\documentclass{article}
\usepackage[UTF8]{ctex}
\usepackage{amsmath}
\usepackage{amssymb}
\usepackage{graphicx}
\usepackage{geometry}
\geometry{a4paper, margin=2cm}

\newcommand{\problemname}[1]{\begin{center}\Large\bfseries #1\end{center}}
\usepackage{enumitem}

\begin{document}
\problemname{求不同高度的大气压值}

蒜头君最近了解到海平面的大气压大约为 $100000$ Pa，还知道了海拔高度和大气压的关系。

在海拔 $3000$ 米内，可以近似地认为每升高 $10$ 米，大气压减小 $100$ Pa，想去登山的蒜头君想尽快地算出山上不同高度对应的大气压值。

\textbf{注意：} 海拔是指地面某个地点高出海平面的垂直距离。

\subsection*{输入格式}
输入仅一行，包括一个实数 $h(100\leh\le2500)$ ，表示山的高度。

\subsection*{输出格式}
输出一个实数，表示对应高度的大气压值。

\subsection*{样例输入}
\begin{verbatim}
100.0
\end{verbatim}

\subsection*{样例输出}
\begin{verbatim}
99000
\end{verbatim}


\end{document}
